\section{Welfare and the Public--Private Resilience Substitution}\label{sec:welfare}

\subsection{Exact surplus accounting}

From utility \eqref{eq:utility}, consumer surplus when both firms survive is
\begin{equation}
CS_D=\frac{1+\gamma}{(2+\gamma)^2}.
\end{equation}
With total producer surplus $2/(2+\gamma)^2$, total surplus in the duopoly state is
\begin{equation}
T_D=\frac{3+\gamma}{(2+\gamma)^2}.
\label{eq:TD}
\end{equation}
When only one firm survives,
\begin{equation}
CS_M=\frac18,\qquad PS_M=\frac14,\qquad T_M=\frac38.
\label{eq:TM}
\end{equation}

Let
\begin{equation}
m_i=s_i(1-r_i)
\end{equation}
be firm $i$'s probability of being unavailable. Under the conditional-independence backup-success law in Section~\ref{sec:model}, the probability that both firms are unavailable is
\begin{equation}
z=J(1-r_1)(1-r_2).
\end{equation}
The four final availability probabilities are therefore $1-m_1-m_2+z$, $m_2-z$, $m_1-z$, and $z$. Expected product-market surplus net of backup costs is
\begin{equation}
\mathcal S
=T_D+(T_M-T_D)(m_1+m_2)+(T_D-2T_M)z
-\frac{k}{2}(r_1^2+r_2^2).
\label{eq:surplus}
\end{equation}
The welfare comparison therefore reflects real output losses, consumer surplus, producer surplus, and real public and private resilience costs. It is not generated by treating fiscal transfers as social costs.

\subsection{Channel decomposition at the canonical witness}

The central welfare logic can be seen at the no-additional-protection point. Consider a symmetric common increase in $a$ at the witness $\gamma=12/25$. Holding firms' backup readiness fixed, the direct reduction in production risk raises coordinated welfare. Allowing backup readiness to respond reverses the sign. The generated decomposition is reported in Table~\ref{tab:channels}.

\begin{table}[t]
\centering
\caption{Marginal welfare channels at the canonical witness}
\label{tab:channels}
\begin{tabular}{lr}\toprule
Component & Marginal effect \\\midrule
Direct protection effect, redundancy fixed & 0.04146 \\
Endogenous redundancy response & -0.04794 \\
Total coordinated effect & -0.00648 \\
Local-government unilateral payoff effect & 0.00193 \\\bottomrule
\end{tabular}

\end{table}

The direct protection effect is approximately $+0.04146$. The endogenous redundancy response contributes approximately $-0.04794$, leaving a coordinated marginal effect near $-0.00648$. At the same point, a jurisdiction's unilateral marginal payoff is approximately $+0.00193$. The signs, rather than the magnitudes, are the economically relevant objects. The direct protection benefit is positive; the induced reduction in private geographic redundancy is sufficiently strong to make the social total negative; and plant attraction keeps the decentralized marginal incentive positive.

\subsection{Why product-market competition matters}

The scarcity-rent wedge $\Delta(\gamma)$ links the industrial-organization block to the resilience decision. A higher $\gamma$ lowers duopoly profit relative to monopoly profit, thereby increasing $\Delta$. Equations \eqref{eq:rC}--\eqref{eq:rD} then imply that equilibrium backup readiness falls with product substitutability on the maintained symmetric interior branch. This does not mean that competition always lowers resilience in richer environments, nor does it assign a general sign to the interaction between product substitutability and the marginal response of readiness to protection. In the maintained baseline model, backup readiness is valuable partly because it determines who remains active after a localized disruption, so the relative value of competitive survival versus sole survival is an essential determinant of the firm's private resilience incentive.
